% !TeX spellcheck = en_US
%----------------------------------------------------------------------------
%
%   This LaTeX template was created by Christian Krieg
%       <christian.krieg@alumni.tuwien.ac.at>
%
% This template is licensed under the following license: Attribution 4.0
% International (CC BY 4.0)
%
%
% You are free to:
%
%   1. Share --- Copy and redistribute the material in any medium or format
%   2. Adapt --- Remix, transform, and build upon the material for any purpose,
%           even commercially.
%
% This license is acceptable for Free Cultural Works.
%
% The licensor cannot revoke these freedoms as long as you follow the license
% terms.
%
% The entire license text is available at:
%   https://creativecommons.org/licenses/by/4.0/legalcode
%
%   May 2021
%
%----------------------------------------------------------------------------
%
\documentclass[%
	a4paper,%
	twoside%
] {article}
%
%----------------------------------------------------------------------------
%
% Title
%
\title{\LaTeX{} review made transparent with \emph{todonotes}, and \emph{git-latexdiff}}
%
%----------------------------------------------------------------------------
%
% Author
%
\author{Christian Krieg <christian@drkrieg.at>}
%
%-----------------------------------------------------------------------------
%
% Use the 'Libertine' font type
%
\usepackage{libertine}
\usepackage[T1]{fontenc}
\usepackage[utf8]{inputenc}
%
%----------------------------------------------------------------------------
%
% Set page margins - original: 2cm
%
\usepackage{geometry}
\geometry{
    left = 2cm,
    right = 2cm,
    top = 2cm,
    bottom = 2cm,
    marginparwidth=1.75cm,
}
%
%----------------------------------------------------------------------------
%
% Set up glossary
%
\usepackage{glossaries}
\newglossaryentry{dvi}{
  name={DVI},
  description={Device-independent},
  text={DVI},
  first={device-independent (DVI)},
}
\newglossaryentry{pdf}{
  name={PDF},
  description={Portable document format},
  text={PDF},
  first={portable document format (PDF)},
}
%\newglossaryentry{}{
%  name={},
%  description={},
%  text={},
%  first={},
%}
%
%----------------------------------------------------------------------------
%
% Set line spacing
%
\usepackage{setspace}
\setstretch{1.5}
%
%----------------------------------------------------------------------------
%
% Set paragraph: No indentation, but include an empty line
%
\usepackage[parfill]{parskip}
%
%----------------------------------------------------------------------------
%
% Commands to set up review comments
%
\usepackage{newtodo}

\newtodo{alice}{purple}
\newtodo{bob}{yellow}
%
%----------------------------------------------------------------------------
%
% Use the cleverref package -- Load this package as the very last!
%
\usepackage{cleveref}
%
%----------------------------------------------------------------------------
%
% Document body
%
\usepackage{lipsum}
\begin{document}

\maketitle

\LaTeX{} provides a great environment for writing scientific texts. With \LaTeX{},
it is very convenient to set mathematical equations, maintain references to a
bibliography, or define and use glossary terms throughout any document. While
\LaTeX{} compiles \LaTeX{} code to a \gls{dvi} file, the pre-dominant use case today
is to compile \LaTeX{} code into a  \gls{pdf} file (see \Cref{fig:compile-pdf}).

\begin{figure}[htp]
    \centering
    \begin{tikzpicture}
        \node (pdflatex) [draw, minimum height=1cm, minimum width=2cm] {pdflatex};
        \draw [<-] (pdflatex.west) -- ++(-1cm,0) node[left]{\LaTeX{} input};
        \draw [->] (pdflatex.east) -- ++(1cm,0) node[right]{\gls{pdf} output};
    \end{tikzpicture}
    \caption{Compiling \LaTeX{} into \protect\gls{pdf} format}
    \label{fig:compile-pdf}
\end{figure}

In order to maintain high quality and to examine the scientific code of conduct,
any scientific text is subject to review. In the pre-dominant use case of
generating \gls{pdf} files from \LaTeX{} files, there are two possibilities:

\begin{enumerate}
    \item{Review the PDF output}
    \item{Review \LaTeX{} input}
\end{enumerate}


Both review options have their advantages and their disadvantages, and it does
not seem to be obvious that one option outperforms the other option.
\Cref{tbl:review-options} summarizes advantages and disadvantages.

\renewcommand{\arraystretch}{1.5}
\begin{table}[htp]
    \centering
    \setstretch{1}
    \caption{Advantages and disadvantages of reviewing \LaTeX{} input or
        \protect\gls{pdf} output}
    \label{tbl:review-options}
%    \begin{tabular}{lm{5cm}m{5cm}}
    \begin{tabular}{ccc}
    \hline
    \textbf{Target} & \textbf{Advantages} & \textbf{Disadvantages} \\
    \hline
    \LaTeX{}
        & \begin{minipage}{5cm}
            \begin{itemize}
                \item{Directly editable}
            \end{itemize}
          \end{minipage}
        & \begin{minipage}{10cm}
            \begin{itemize}
                \item{Appearance of end product is not visible (e.g., figures)}
                \item{If edited directly, it is not clear what has been changed
                     by the reviewer}
            \end{itemize}
          \end{minipage}
        \\[2em]
    \gls{pdf}
        & \begin{minipage}{5cm}
            \begin{itemize}
                \item{Comments can be made directly in the end product}
                \item{Nearly every PDF viewer supports to include annotations}
            \end{itemize}
          \end{minipage}
        & \begin{minipage}{10cm}
            \begin{itemize}
                \item{\gls{pdf} annotations don't provide an overview over
                    all comments}
                \item{It is hardly possible to reply to comments}
                \item{It is hard to maintain reviews of one document by multiple
                    authors}
                \item{It is nearly impossible to track annotations in revision
                    control systems (because \gls{pdf} files are binary)}
                \item{It is cumbersome to fix minor issues directly (e.g.,
                    wording, commas, etc.). Instead, a review annotation must be
                    given}
            \end{itemize}
          \end{minipage}
        \\
    \hline
    \end{tabular}
\end{table}

 As a
plain text format, it is very convenient to track \LaTeX{} documents with revision
control systems


\end{document}
%
%----------------------------------------------------------------------------
